% !TEX root = ../main.tex
\chapter{\ifenglish Conclusions and Discussions\else บทสรุปและข้อเสนอแนะ\fi}

ในบทนี้สรุปผลการดำเนินงานของโครงงานระบบการสื่อสารความปลอดภัยสูงสำหรับการลงคะแนนเสียงโดยใช้การกระจายกุญแจแบบควอนตัม (BB84 QKD) ทั้งในด้านความสำเร็จตามวัตถุประสงค์ ผลการทดสอบบนฮาร์ดแวร์จริงของ IBM Quantum ปัญหาอุปสรรคพร้อมแนวทางแก้ไขเชิงวิศวกรรม ข้อจำกัด และข้อเสนอแนะสำหรับการพัฒนาต่อยอด

\section{\ifenglish Conclusions\else สรุปผลการดำเนินงาน\fi}

โครงงานได้ออกแบบ พัฒนา และทดสอบระบบการลงคะแนนเสียงอิเล็กทรอนิกส์ที่มีปลอดภัยแบบ Information-Theoretic Security โดยนำโพรโทคอล BB84 ร่วมกับการเข้ารหัส One-Time Pad (OTP) มาประยุกต์ใช้ในช่องทางการส่งบัตรลงคะแนน โดยมีผลสัมฤทธิ์หลักตามวัตถุประสงค์ดังนี้:

\begin{enumerate}
    \item \textbf{กระบวนการความปลอดภัยของระบบ:} พัฒนาระบบตั้งแต่การสร้างวงจรควอนตัมใน Z-basis และ X-basis, การเปรียบเทียบแกน (Basis Sifting), การประเมินค่า QBER จากบิตทดสอบ, การตรวจจับผู้ดักฟัง, การขยายความปลอดภัยด้วยเมทริกซ์โทพลิตซ์ (Privacy Amplification) และการเข้ารหัสบัตรลงคะแนนด้วย One-Time Pad ($C = M \oplus K$) ได้อย่างสมบูรณ์
    
    \item \textbf{การทดสอบบนฮาร์ดแวร์ IBM Quantum จริง:} ผลการทดสอบบนชิปประมวลผลควอนตัมยืนยันประสิทธิภาพของระบบ 3 ประการ:
    \begin{itemize}
        \item \textit{การตรวจจับการดักฟัง:} ตรวจจับการดักฟังแบบ Intercept-and-Resend ได้แม่นยำ $100.0\%$ ปราศจาก False Alarm เมื่อขนาดกุญแจ $n \ge 32$ บิต
        \item \textit{ความเที่ยงตรงของบิตกลุ่มตัวอย่าง:} ค่าความคลาดเคลื่อนเฉลี่ยในการประมาณการ QBER ลดลงตามกฎจำนวนมาก จาก $16.69\%$ ที่ $n = 8$ บิต เหลือเพียง $2.79\%$ ที่ $n = 256$ บิต
        \item \textit{การกำจัดข้อมูลรั่วไหล:} กระบวนการ Privacy Amplification ด้วยค่าความปลอดภัย $s \ge 16$ บิต สามารถลดข้อมูลที่ผู้ดักฟังล่วงรู้ได้เกือบทั้งหมด (ความสำเร็จ $99\% - 100\%$)
    \end{itemize}

    \item \textbf{เว็บแอปพลิเคชันเชิงโต้ตอบ:} พัฒนาระบบส่วนต่อประสาน 3 มุมมอง (Voter View, Practical View, และ System God View) ช่วยให้สังเกตและตรวจสอบการทำงานของวงจรควอนตัมและสถานะความปลอดภัยได้อย่างชัดเจน
\end{enumerate}

\subsection{ข้อจำกัดของระบบที่ค้นพบจากการทดลอง}

\begin{itemize}
    \item \textbf{ความเป็นนิรนาม (Anonymity Trade-off):} การถอดรหัสแบบ One-Time Pad ฝั่งผู้รับ (Server) จำเป็นต้องจับคู่กุญแจกับผู้ลงคะแนน ($M = C \oplus K$) เซิร์ฟเวอร์จึงยังล่วงรู้ความสัมพันธ์ระหว่างผู้ลงคะแนนกับบัตรเลือกตั้งได้ ซึ่งต้องใช้กลไกไม่เปิดเผยตัวตนในระดับแอปพลิเคชันเพิ่มเติมหากต้องการความเป็นนิรนามขั้นสูงสุด
    \item \textbf{การไม่มีกลไกแก้ไขข้อผิดพลาด (No Error Correction):} ระบบอาศัยการยกเลิกเซสชัน (Abort) ทันทีที่ QBER เกินเพดาน $11.0\%$ โดยไม่มีโพรโทคอลแก้ไขข้อผิดพลาดทางคลาสสิก (เช่น Cascade หรือ LDPC) ทำให้ประสิทธิภาพขึ้นอยู่กับคุณภาพของชิปควอนตัมโดยตรง
    \item \textbf{ระยะทางของช่องสัญญาณกายภาพ:} ในสายใยแก้วนำแสงจริง อนุภาคโฟตอนเดี่ยวจะเกิดการลดทอนสัญญาณตามระยะทาง (Attenuation) จำกัดระยะการเชื่อมต่อระหว่างหน่วยเลือกตั้งกับศูนย์กลางไว้ที่ไม่เกิน 50--100 กิโลเมตร
\end{itemize}

\section{\ifenglish Challenges and Engineering Solutions\else ปัญหาที่พบและแนวทางการแก้ไขเชิงวิศวกรรม\fi}

ในการพัฒนาและทดสอบระบบบนฮาร์ดแวร์ประมวลผลควอนตัมจริง พบความท้าทายทางวิศวกรรมและมีแนวทางแก้ไข 4 ประการ:

\begin{enumerate}
    \item \textbf{การจำแนกสัญญาณรบกวนฮาร์ดแวร์ออกจากการดักฟัง:} สัญญาณรบกวนทางกายภาพของฮาร์ดแวร์จริง ($2.0\% - 3.0\%$) อาจกระตุ้นให้ระบบแจ้งเตือนเท็จ โครงงานแก้ปัญหานี้โดยกำหนดเกณฑ์เพดาน $QBER_{\text{threshold}} = 11.0\%$ ตามบทพิสูจน์ของ Shor และ Preskill ซึ่งสร้างส่วนต่างความปลอดภัย (Safety Margin) กว้าง $\approx 8.5\%$ เหนือสัญญาณรบกวน ทำให้สามารถจำแนกการดักฟังแบบ Intercept-and-Resend ($\approx 25\%$) ได้แม่นยำ $100\%$ โดยไม่มีการแจ้งเตือนเท็จ
    
    \item \textbf{ข้อจำกัดจำนวนคิวบิตและเวลารอคิวบนคลาวด์ควอนตัม:} การสร้างกุญแจยาวในรอบเดียวทำได้ยากเนื่องจากข้อจำกัดคิวบิตและเวลารอคิว จึงพัฒนาเทคนิค \textbf{Multi-Session Key Concatenation} โดยแบ่งการสร้างกุญแจออกเป็นรอบย่อยขนาด 20--40 คิวบิต แล้วนำกุญแจที่ผ่านการตรวจสอบมาเชื่อมต่อกันจนครบความยาวเป้าหมาย ซึ่งใช้รอบย่อยเฉลี่ยเพียง $3.2$ รอบสำหรับกุญแจ 24 บิต
    
    \item \textbf{ความแปรปรวนทางสถิติในกุญแจขนาดเล็ก:} ที่ขนาดกุญแจสั้น ($n = 8$ บิต) บิตกลุ่มตัวอย่างมีน้อย (9 บิต) ทำให้เกิดความผันผวนทางสถิติจนเกิด False Alarm $7\%$ และมีอัตราหลุดรอด $10\%$ จึงกำหนดข้อกำหนดทางวิศวกรรมให้ $n \ge 32$ บิตขึ้นไป (กลุ่มตัวอย่าง $\ge 38$ บิต) ซึ่งลดความคลาดเคลื่อนเฉลี่ยต่ำกว่า $7.8\%$ และขจัดความผิดพลาดทั้งสองส่วนเหลือ $0.0\%$
    
    \item \textbf{การรั่วไหลของข้อมูลในการดักฟังแบบแอบแฝง (Sub-Threshold):} หากผู้ดักฟังดักข้อมูลเพียงบางส่วนเพื่อให้ QBER ไม่เกิน $11.0\%$ ผู้ดักฟังอาจได้ข้อมูลกุญแจดิบไปบางส่วน ($4.8$ ถึง $51.6$ บิต) โครงงานแก้ปัญหาโดยใช้ Toeplitz Universal Hashing บน $\text{GF}(2)$ ด้วยค่าพารามิเตอร์ $s \ge 16$ บิต ซึ่งบีบอัดและทำลายข้อมูลรั่วไหลจนผู้ดักฟังมีความรู้หลังการบีบอัดเป็น $0.00$ บิตอย่างแน่นอน
\end{enumerate}

\section{\ifenglish Suggestions and Future Work\else ข้อเสนอแนะและแนวทางการพัฒนาต่อ\fi}

แนวทางการวิจัยและพัฒนาต่อยอดเพื่อยกระดับสู่การใช้งานจริงในอนาคต ประกอบด้วย 4 มิติ:

\begin{enumerate}
    \item \textbf{การยกระดับความเป็นนิรนามขั้นสูง:} ประยุกต์ใช้เทคนิคทางวิทยาการรหัสลับสมัยใหม่ เช่น เครือข่ายสับเปลี่ยนบัตรลงคะแนน (Mix-nets) หรือการพิสูจน์แบบความรู้เป็นศูนย์ (Zero-Knowledge Proofs: ZKP) เพื่อตัดความสัมพันธ์ระหว่างตัวตนผู้ลงคะแนนกับบัตรเลือกตั้งในระดับแอปพลิเคชัน
    
    \item \textbf{การผสานอัลกอริทึมแก้ไขข้อผิดพลาดทางคลาสสิก:} ผสานรวม Cascade Protocol หรือ Low-Density Parity-Check (LDPC) Codes ในขั้นตอน Classical Post-Processing เพื่อกู้คืนบิตที่ผิดพลาดจากสัญญาณรบกวนโดยไม่ต้องทำลายกุญแจทิ้ง ช่วยเพิ่มอัตราการสร้างกุญแจสำเร็จ (Key Generation Rate)
    
    \item \textbf{การทดสอบบนโครงข่ายใยแก้วนำแสงจริง:} นำระบบไปเชื่อมต่อกับฮาร์ดแวร์โฟโตนิกส์จริง ได้แก่ เลเซอร์พัลส์อ่อน ตัวปรับมุมโพลาไรเซชัน สายไฟเบอร์ออปติก Single-Mode และตัวตรวจวัดโฟตอนเดี่ยว (SPADs) เพื่อประเมินผลกระทบของการลดทอนสัญญาณจริง
    
    \item \textbf{สถาปัตยกรรมความปลอดภัยแบบลูกผสม (Hybrid QKD + PQC):} นำวิทยาการรหัสลับยุคหลังควอนตัมตามมาตรฐาน NIST เช่น ML-DSA (Dilithium) มาใช้ยืนยันตัวตนคู่สื่อสารบนช่องทางคลาสสิก เพื่อป้องกันการโจมตีแบบสวมรอยเป็นคนกลาง (Man-in-the-Middle) และใช้ ML-KEM (Kyber) ร่วมกับ QKD เพื่อสร้างการป้องกันแบบหลายชั้น (Defense-in-Depth)
\end{enumerate}
